% Source-derived chapter 06: Families of Enriques surfaces
% Original source: no-enriques-surface-over-the-integers
\chapter{Families of Enriques surfaces}
\label{no-enriques-surface-over-the-integers-chapter-06}
\source{no-enriques-surface-over-the-integers}

\begin{remark}[Source section]
\label{no-enriques-surface-over-the-integers-chapter-06-source}
\source{no-enriques-surface-over-the-integers}
\source{no-enriques-surface-over-the-integers}
This chapter reproduces the corresponding section of the retrieved
LaTeX source, including its definitions, statements, examples, and
proofs. It has not yet been translated into Lean.
\notready
\end{remark}

% The source section title was: Families of Enriques surfaces
\mylabel{Families}

In this section we examine    arithmetic properties of Enriques surfaces $Y$
over non-closed ground fields. In fact, we treat them as a special case of the following general notion:
A \emph{family of Enriques surfaces} over an arbitrary ground ring $R$ is an algebraic space $\foY$,
together with a morphism  $f:\foY\ra\Spec(R)$
that is flat,  proper and of finite presentation such that for each geometric point $\Spec(\Omega)\ra\Spec(R)$,
the base-change $Y$ is an Enriques surfaces over $k=\Omega$.  
Let $\shM_\Enr$ be the resulting stack in groupoids 
whose fiber categories $\shM_\Enr(R)$ consists 
of the families of Enriques surfaces $f:\foY\ra\Spec(R)$.  This stack lies over the site $(\Aff/\ZZ)$ of
all affine schemes, endowed with the fppf topology. 
We already   can formulate the main result of this paper:

\begin{theorem}
\source{no-enriques-surface-over-the-integers}
\notready
\mylabel{no enriques over integers}
The fiber category $\shM_\Enr(\ZZ)$ is empty. In other words, 
there is no family of Enriques surfaces over the ring $R=\ZZ$.
\end{theorem}

The proof is long and indirect, and will  finally be achieved  in Section \ref{Proof}. 
The following result  (\cite{Schroeer 2021b}, Theorem 7.2) tells us that we do not have to worry
about 
exceptional Enriques surface:

\begin{theorem}
\source{no-enriques-surface-over-the-integers}
\notready
\mylabel{no exceptional}
There is no family of Enriques surfaces over the ring $R=\ZZ/4\ZZ$ whose
geometric fiber is exceptional.
\end{theorem}

We now collect further preliminary facts.
Throughout, $f:\foY\ra \Spec(R)$ denotes a family of Enriques surfaces over an arbitrary ground ring $R$.
To simplify notation, we set $S=\Spec(R)$.
According to \cite{Artin 1969}, Theorem 7.3 
the fppf sheafification of the naive Picard functor $A\mapsto\Pic(\foY\otimes_RA)$ is representable by a 
group object in the category of algebraic spaces over $R$.  
The following observation, which immediately follows from \cite{Ekedahl; Hyland; Shepherd-Barron 2012}, Corollary 4.3,
will be important:

\begin{proposition}
\source{no-enriques-surface-over-the-integers}
\notready
The algebraic space $\Pic_{\foY/R}$ is a scheme.
The sheaf of subgroups $\Pic^\tau_{\foY/R}$ is representable by an open embedding,
and its structure morphism  is  locally free of rank  two. 
Moreover, the     quotient sheaf $\Num_{\foY/R}$ is representable by a local system of free abelian groups of rank $\rho=10$.
\end{proposition}
 
Each triple $(L,a,b)$ where $L$ is an invertible $R$-module and $a\in L$, $b\in L^{\otimes -1}$ are
elements with $a\otimes b=2_R$ yields a finite flat group scheme $G^L_{a,b}$, characterized by
$$
G^L_{a,b}(A)=\{f\in L\otimes_RA\mid f^{\otimes 2}=a\otimes f \},
$$
with group law $f_1\star f_2=f_1+f_2+ b\otimes  f_1\otimes   f_2$.  
According to \cite{Tate; Oort 1970}, Theorem 2 each 
group scheme that is locally free of rank two is isomorphic to some $G^L_{a,b}$. Cartier duality corresponds
to the involution $(L,a,b)\mapsto (L^{\otimes -1},b,a)$.
The special case $L=R$, $a=1$, $b=2$ yields the constant group scheme $(\ZZ/2\ZZ)_R$,
whereas $a=2$, $b=1$ defines the multiplicative group scheme $\mu_{2,R}$.

Write $P=\Pic^\tau_{\foY/R}$,   and   $G=\uHom(P,\GG_{m,R})$ for the Cartier dual.
According to \cite{Raynaud 1970}, Proposition 6.2.1, the inclusion $P\subset\Pic_{\foY/R}$ corresponds to a global section
of $R^1f_*(G_\foY)$, which we denote by a formal symbol $[\foX]$. This  notation may  be explained as follows:
 The Leray--Serre spectral sequence for the structure morphism $f:\foY\ra S=\Spec(R)$
yields an exact sequence
\begin{equation}
\label{5-term sequence}
\source{no-enriques-surface-over-the-integers}
0\ra H^1(S, G)\ra H^1(\foY,G_\foY)\ra H^0(S,R^1f_*(G_\foY))\stackrel{d}{\ra} H^2(S,G)\ra H^2(\foY,G_\foY).
\end{equation}
If $[\foX]\in H^0(S,R^1f_*(G_\foY))$ maps to zero under the differential $d$, the section comes
from a $G_\foY$-torsor $\epsilon:\foX\ra \foY$, which we call a \emph{family of canonical coverings}. 

\begin{proposition}
\source{no-enriques-surface-over-the-integers}
\notready
\mylabel{canonical covering}
A family of canonical coverings $\epsilon:\foX\ra\foY$ exists 
under any of the following three conditions:
\begin{enumerate}
\item The pullback map $H^2(S,G)\ra H^2(\foY,G_\foY)$ is injective.
\item The   structure morphism $f:\foY\ra S$ admits a section.
\item The group scheme $P\ra S$ is \'etale and $\Pic(S)=0$.
\end{enumerate}
If a family of canonical coverings exists,  it is unique up to isomorphism and twisting by  pullbacks of $G$-torsors over $S$.
\end{proposition}

\proof
The first assertion and the statement about uniqueness immediately follow from
the five-term exact sequence \eqref{5-term sequence}.
If the structure morphism admits a section, the map $H^2(S,G)\ra H^2(\foY,G_\foY)$ has a retraction and is thus injective,
which gives (ii).

Finally, suppose that $P\ra S$ is \'etale. Suppose first that $R$ is noetherian.
Consider the invertible sheaf $\shL=\Omega^2_{\foY/R}$.
For each $a\in S$, the restriction $\shL|Y$ to the fiber $Y=f^{-1}(a)$ is the dualizing sheaf, which has order two
in $\Pic(Y)$. Moreover, we have $h^i(\shL^{\otimes 2}|Y)=h^i(\O_Y)=0$ for all $i\geq 1$. In turn, the formation of the direct image
$\shN=f_*(\shL^{\otimes 2})$ commutes with base-change, and we conclude that $\shN$ is invertible.
With Nakayama, we infer that $f^*(\shN)\ra \shL^{\otimes 2}$ is bijective.
Using $\Pic(R)=0$ we obtain a trivialization $\varphi:\O_\foY\ra\shL^{\otimes 2}$.
We endow the locally free sheaf
$\shA=\O_\foY\oplus\shL^{\otimes -1}$
with an  algebra structure, by declaring $(a,b)\cdot (a',b') = (aa'+\varphi^{-1}(bb'), ab'+a'b)$.
Then the relative spectrum $\foX=\Spec(\shA)$ is the desired canonical covering.

It remains to treat general rings $R$. Since   $f:\foX\ra\Spec(R)$ is of finite presentation, it comes
from a family of Enriques surfaces over some noetherian subring $R'$. We then can argue as in the previous paragraph,
after enlarging the noetherian ring $R'$ to make $\shN$ trivial. 
\qed

\iffalse
Under suitable assumptions, one can say more. Here the Picard group $\Pic(R)=H^1(S,\GG_m)$ 
and the Brauer group $\Br(R)=H^2(S,\GG_m)$ become highly relevant. 
As customary, we write   $M[2]$ and $M/2M$ 
for kernel and  cokernel for the multiplication-by-two map on an abelian group $M$.

\begin{proposition}
\source{no-enriques-surface-over-the-integers}
\notready
Suppose that we have
$$
\Pic^\tau_{\foY/R}=\mu_{2,R}\quadand \Pic(R)/2\Pic(R)=0\quadand   \Br(R)[2]=0.
$$
Then  the set of isomorphism classes 
of canonical coverings $\epsilon:\foX\ra\foY$ is a principal homogeneous space for the group $R^\times/R^{2\times}$.
In particular, such a covering exists. 
\end{proposition}

\proof
\qed
\fi

\medskip
We say that our family of Enriques surfaces has \emph{constant local system} if 
the local system $\Num_{\foY/R} $ is isomorphic  to  
$(\ZZ^{\oplus 10})_R$. Likewise, we say that it has \emph{constant Picard scheme} if 
the group space $\Pic_{\foY/R}$ is isomorphic to $(\ZZ/2\ZZ\times\ZZ^{\oplus 10})_R$.
Moreover, we say that it has \emph{split Picard scheme} if the short exact sequence of group spaces
\begin{equation}
0\lra \Pic^\tau_{\foY/R}\lra \Pic_{\foY/R}\lra \Num_{\foY/R}\lra 0
\end{equation}
splits. Our interest in these notions stems from the following fact: 

\begin{proposition}
\source{no-enriques-surface-over-the-integers}
\notready
\mylabel{constant picard scheme}
If the base ring is $R=\ZZ$, then our family of Enriques surfaces $f:\foY\ra\Spec(R)$ has constant Picard scheme.
\end{proposition}

\proof
Choose  a geometric point $\Spec(\Omega)\ra S$.
The local system $\Num_{\foY/R}$ comes from a representation of the algebraic fundamental group
$\pi_1(S,\Omega)$ on the stalk $\Num_{Y/k}(\Omega)=\ZZ^{\oplus 10}$. According to Minkowski's theorem  
(see for example \cite{Neukirch 1999}, Theorem 2.18), the scheme $S=\Spec(\ZZ)$
is simply-connected, so our family of Enriques surfaces has constant local system.

For the ring $R=\ZZ$, up to isomorphism the possible triples describing the group scheme $P=\Pic^\tau_{\foY/R}$
are either $(R,1,2)$ or $(R,2,1)$. In other words, we either have $P=(\ZZ/2\ZZ)_R$ or $P=\mu_{2,R}$.
Seeking a contradiction, we assume the latter holds.
The Picard group $\Pic(R)$ vanishes, because the ring  $R=\ZZ$ is a principal ideal domain.
Furthermore, we have $\Br(R)=0$. Indeed, the Hasse principle (see for example \cite{Neukirch; Schmidt; Wingberg 2008}, Theorem 8.1.17)
gives a short exact sequence
$$
0\lra\Br(\QQ)\lra \bigoplus\Br(\QQ_\primid)\stackrel{\delta}{\lra} \QQ/\ZZ\lra 0,
$$
where the sum runs over all places $\primid$ including the Archimedean one. Since we have $\Br(\ZZ_p)=\Br(\FF_p)=0$ for all primes $p>0$ by Wedderburn,
one sees that $\Br(\ZZ)$ is contained in the Brauer group $\Br(\RR)\cap\Kernel(\delta)=0$.
 
Over the localization $R'=R[1/2]$, the group scheme $\mu_{2,R}$ becomes \'etale, hence there is a canonical covering
$\foX\otimes R'\ra\foY\otimes R'$, by Proposition \ref{canonical covering}.
Now consider the generic fiber $X=\foX\otimes\QQ$, which is a K3 surface. In particular, the Hodge number
$h^{2,0}(X)=h^0(\Omega_X^2)$ is non-zero.
By construction, $X$ has good reduction over the localization $\ZZ_{(p)}$ for all primes $p\geq 3$.
We now consider the ring of Witt vectors 
$A=W(\FF_2^\sep)$, which is the maximal unramified extension of the local ring $\ZZ_2$.

Recall that Fontaine showed that for each proper smooth scheme over $\QQ$ that has good reduction over
$W(\FF_p^\sep)$ for all primes $p>0$, the Hodge numbers $h^{i,j}$ must vanish for $i\neq j$ and $i+j\leq 3$
(\cite{Fontaine 1993}, Theorem 1 on page 44, and first remark afterwards).
Essentially the same  result was independently obtained by Abrashkin, who showed that for each smooth proper
scheme over $\ZZ$, the Hodge numbers $h^{i,j}$ of  the complex fiber vanish for $i\neq j$ and $i+j\leq 3$
(\cite{Abrashkin 1990}, \S7,  Section 6, Theorem on page 516. Note that the formulation given there contains an obvious misprint,
confer the Theorem on page 514).

Since our K3 surface $X=\foX\otimes\QQ$ has Hodge number $h^{2,0}=1$ and   good reduction over   $W(\FF_p^\sep)$ for all $p\neq 2$, it follows that 
the base-change $X_F$ to the field of fractions $F=\Frac(A)$ 
must have bad reduction over $A$. 

One the other hand, using Hensel's Lemma, any $\FF_2^\alg$-valued point for $\foY\ra S$ can be extended to an $A$-valued point.
Again with Proposition \ref{canonical covering} we conclude that a family of canonical covering  $\foX\otimes A\ra\foY\otimes A$ exists,
whence $X_F$ has good reduction over $A$, contradiction.

Summing up, the Picard scheme sits in  a short exact sequence
\begin{equation}
\label{extension local systems}
\source{no-enriques-surface-over-the-integers}
0\lra (\ZZ/2\ZZ)_R\lra \Pic_{\foY/R}\stackrel{\pr}{\lra} (\ZZ^{\oplus 10})_R\lra 0.
\end{equation}
For each global section $l$ on the right, the preimage $T=\pr^{-1}(l)$ is a representable torsor for
the group scheme $H=(\ZZ/2\ZZ)_R$. In particular, the structure morphism $T\ra \Spec(R)$ is finite
\'etale, of degree two. Since $\pi_1(S,\Omega)=0$, the total space $T$ is the disjoint union of two copies
of $S=\Spec(\ZZ)$. From this we infer that the short exact sequence \eqref{extension local systems} splits.
In turn, our family of Enriques surfaces has constant Picard scheme.
\qed

\medskip
When combined with suitable assumptions
on the Picard group or the Picard group of the base ring $R$, our constancy condition
on   the Picard scheme have remarkable consequence for the fibers of the
morphism $f:\foY\ra\Spec(R)$.
In what follows, let $R\ra k$ be a homomorphism
to a field, choose   an algebraic closure $\bk=k^\alg$,
and write  $Y=\foY\otimes_Rk$  and $\bY=\foY\otimes_R \bk$ for the resulting Enriques surfaces.

\iffalse
\begin{proposition}
\source{no-enriques-surface-over-the-integers}
\notready
Suppose that the  groups $\Pic(R)/2\Pic(R)$ and $\Br(R)[2]$ vanish.
Assume furthermore that our family of Enriques surfaces $f:\foY\ra\Spec(R)$ either has constant local system
and $\Pic^\tau_{\foY/R}=\mu_{2,R}$, or constant Picard scheme.
Then the set of rational points $Y(k)$ is non-empty.
\end{proposition}

\begin{corollary}
\source{no-enriques-surface-over-the-integers}
\notready
Assumptions as in the Proposition. If $R$ is a Dedekind ring, then the structure morphism
$f:\foY\ra\Spec(R) $ admits a section.
\end{corollary}
\fi

\begin{theorem}
\source{no-enriques-surface-over-the-integers}
\notready
\mylabel{family with constant pic}
Suppose that the family of Enriques surfaces  $f:\foY\ra\Spec(R)$ has constant Picard scheme.
Assume furthermore  that  $\Pic(R)$  and  $\Br(R)$ vanish.
Then there is at least one canonical covering $\epsilon:\foX\ra \foY$, and 
the set of isomorphism classes  of such coverings  
is a principal homogeneous space for the  group
$R^\times/R^{\times 2}$. Furthermore, the induced Enriques surfaces $Y$ and $\bY$ enjoy the following properties:
\begin{enumerate}
\item 
The dualizing sheaf $\omega_Y$   has order two in the Picard group.
\item 
Every class $l\in \Num_{Y/k}(\bk)$ comes from an invertible sheaf $\shL$ on $Y$,
which is unique up to twisting by $\omega_Y$.
\item 
Every $(-2)$-curve $\bE\subset\bY$ is the base-change of a $(-2)$-curve $E\subset Y$, which is isomorphic  
to the projective line $\PP^1_k$.
\item 
Every genus-one fibration $ \bY\ra\PP^1_{\bk}$ is the base-change of a genus-one fibration
$ Y\ra\PP^1_k$. The latter has exactly two multiple fibers, both of which are tame and lie over   $k$-rational points.
\item 
There is at least one genus-one fibration $\varphi:Y\ra\PP^1$.
\end{enumerate}
\end{theorem}

\proof
By assumption we have $\Pic^\tau_{\foY/R}=(\ZZ/2\ZZ)_R$. 
Write $S=\Spec(R)$. The Kummer sequence $0\ra \mu_2\ra\GG_m\stackrel{2}\ra\GG_m\ra 0$ gives an exact sequence
$$
\Pic(S)\stackrel{2}{\lra}\Pic(S)\lra H^2(S,\mu_{2})\lra \Br(S)\stackrel{2}{\lra}\Br(S).
$$
Our assumptions ensure that the term in the middle vanishes.
According to Proposition \ref{canonical covering},   there is at least one canonical covering $\epsilon:\foX\ra\foY$.
Moreover, the set of isomorphism classes is a principal homogeneous space for the group $H^1(S,\mu_2)$.
Again by the Kummer sequence we obtain an identification $H^1(S,\mu_2)=R^\times/R^{\times 2}$.
This establishes the statements on the family of canonical coverings.

From $\Pic^\tau_{\foY/R}=(\ZZ/2\ZZ)_R$ we immediately get (i).
To see (ii), we may regard the element $l$ as a section of $R^1f_*(\GG_{m,\foY})$.
The Leray--Serre spectral sequence for the structure morphism $f:\foY\ra S$ yields
$$
H^1(\foY,\GG_{m,\foY})\lra H^0(S, R^1f_*(\GG_{m,\foY}))\lra H^2(S,\GG_{m,S}).
$$
The arrow on the right is the zero map  by assumption, hence the section $l$ comes from an invertible sheaf on $\foY$.
Base-changing along $R\ra k$, we find the desired invertible sheaf $\shL$ on $Y$.

% We next check that $Y$ contains a rational point. Suppose first that there is a $(-2)$-curve $E'\subset \bY$.
% Since the intersection form on $\Num(\bY)$ is unimodular, there is an invertible sheaf $\shL'$ on $\bY$
% with $(\shL'\cdot E')=1$.

We now check (iii).  Let $\bE\subset\bY$ be a $(-2)$-curve, and consider the invertible sheaf $\bar{\shL}=\O_{\bY}(\bE)$.
We just saw that it is  the base-change of some invertible sheaf $\shL$ on $Y$.
Using $h^0(\bar{\shL})=1$ we infer that there is a unique curve $E\subset Y$ inducing $\bE\subset\bY$, and therefore $E^2=-2$.
This curve is  a twisted form of the projective line. Since its class in $\Num(Y)$ must be primitive in light of the self-intersection number,
and  the intersection form is unimodular,   there
is an invertible sheaf $\shN$ with $(\shN\cdot E)=1$, hence $E\simeq \PP^1$.

It remains to verify (iv). 
Let $\bY\ra\PP^1_{\bar{k}}$ be a genus-one fibration.
According to Theorem \ref{projective contractions}, it is the base-change of a fibration $Y\ra B$, where $B$ is a twisted form
of the projective line.  Consider the invertible sheaf  $\bar{\shL}=\O_\bY(\bar{F})$, where $\bar{F}\subset\bY$
is a half-fiber. It arises from some invertible sheaf $\shL$ on $Y$.   Using $h^0(\bar{\shL})=1$ we infer that there is a unique curve
$F\subset Y$ inducing $\bar{F}\subset \bY$. This curve is a half-fiber for $Y\ra B$. 
If $\bar{F}$ is reducible, it contains a $(-2)$-curve. By the previous paragraph, this gives an inclusion $\PP^1\subset F$.
In particular, $Y$ contains a rational point, hence $B=\PP^1$. Now suppose that $\bar{F}$ is irreducible.
Then $F$ is an integral curve with $h^0(\O_F)=h^1(\O_F)=1$. Since $\Num(Y)$ is unimodular, there
is an invertible sheaf $\shN$ on $Y$ with $(\shN\cdot F)=1$. With Riemann--Roch we conclude that there
is an effective Cartier divisor of degree one on $F$. Again there is a rational point, and $B=\PP^1$.
Summing up, we have shown that each half-fiber on $\bY$ induces a half-fiber on $Y$ that lies over a rational point.
Since there are exactly two half-fibers on $\bY$, both of which are tame, the same holds for $Y$. 
\qed

\iffalse
\medskip
Except for the computation of $H^1(S,G)$, the above arguments also work  under the weaker but more technical assumptions
that $\Pic(R)$ is 2-divisible, that $\Br(R)$ is 2-torsion-free, and that the pullback map $\Br(S)\ra\Br(\foY)$
is injective.
Also note that the our result applies, in particular, for the ring of integers $R=\ZZ$ and for local henselian rings
with $\Br(R/\maxid_R)=0$,   in particular for finite fields $R=\FF_q$.
\fi


%===========================================================
